% !TeX root = ../../Book5.tex
% 纲要标题：## 附录 D　微分算子、D-模与几何解析方法
% 纲要位置：工作记录/计划纲要.md，第 5024 行。

\chapter{微分算子、D-模与几何解析方法}
\label{b5:app:D}
\BookChapterNumber{b5:app:D}

\begin{chapterabstract}
本附录以微分算子模研究线性微分方程的代数约束，计算阶滤过与特征簇，再比较形式、解析及渐近解，并介绍 Stokes 现象与正则情形的对应。
\end{chapterabstract}

\begin{learninggoals}
\begin{itemize}
\item 会从线性方程或方程组构造左微分算子模，并用同态恢复解。
\item 区分阶滤过与 Bernstein 滤过，计算主理想商模的特征簇。
\item 根据特征簇检验 holonomic 条件，辨别联络、点支撑模及一般凝聚 D-模。
\item 比较形式存在性、解析单值和扇形渐近数据，用非正则算例说明正则条件的作用。
\end{itemize}
\end{learninggoals}

在解析函数中，\(y'=0\) 有常数解，而 \(xy=0\) 只有零解；若允许点支撑对象，后一方程又会出现别的行为。把两条方程写成微分算子环 \(D\) 上的左模 \(D/D\partial\) 与 \(D/Dx\)，可以在指定解空间以前先研究算子之间的代数关系。主符号和特征簇将给出第一层几何信息；随后还要说明，局部化、解析延拓与渐近条件怎样改变实际能看到的解。

\ProseParagraphBreak
代数部分需要左模、滤过与伴随分次、仿射簇和层的基本语言；最后的解析讨论另需复函数的 Laurent 展开、解析延拓及留数定理。导出 Hom 算例只需自由消解，一般对应的概要则涉及消解奇点、GAGA 和导出层论。

% !TeX root = ../../Book5.tex
% 纲要标题：### D.1 从微分模到 D-模
% 纲要位置：工作记录/计划纲要.md，第 5026 行。
% 纲要细目（保留原定写作范围）：
%
% * Weyl 代数模、微分算子环及线性方程组；用点支撑分布比较解对象，说明同时零化位置与微分方向的限制
% * 代数微分算子、滤过与伴随分次的符号思想
% * 特征零光滑代数簇上的 D-模与特征簇；明确与线性方程组的联系
% * holonomic D-模的定义和代表例子；以二维偏导系统比较特征簇维数。holonomic 保留原文，不译作 holomorphic（全纯）；也不是第三卷微分代数的同义词
% #### D.1.1 Weyl 代数模与线性方程组
% #### D.1.2 微分算子、滤过与伴随分次
% #### D.1.3 光滑代数簇上的 D-模与特征簇
% #### D.1.4 holonomic D-模与术语辨析
%
\section{从微分模到 D-模}
\label{b5:sec:D01}

线性方程 \(P(y)=0\) 中，未知量所在的函数空间可以改变，算子之间的代数关系却可以预先记录下来。D-模把这些关系组成一个模，再通过模同态求解。以下先在仿射直线上计算，随后说明在光滑簇上需要添加哪些几何结构。

\subsection{Weyl 代数模与线性方程组}
\label{b5:subsec:D01:01}

设 \(k\) 为特征零域。在 \(k[x]\) 上，把乘以 \(x\) 的算子仍记为 \(x\)，把求导算子记为 \(\partial\)。Leibniz 法则给出
\[
 \partial x-x\partial=1.
\]
这些算子生成的 \(k\) 代数为一阶 \term[Weyl-daishu]{Weyl 代数}[Weyl algebra] \(A_1(k)\)。它的每个元素唯一写为
\(\sum_{j=0}^m a_j(x)\partial^j\)。存在性来自交换关系不断把 \(x\) 移到左边。若这种表达为零算子，取最大 \(m\)，对该算子反复与乘法算子 \(x\) 取交换子 \(m\) 次，便得乘法算子 \(m!\cdot a_m(x)\) 为零，从而 \(a_m=0\)，归纳得所有系数为零。这也证明抽象生成关系没有额外隐藏关系。

\begin{proposition}\label{b5:prop:D-solution-module}
设 \(k\) 为特征零域，\(D=A_1(k)\)，\(P\in D\)，\(V\) 为左 \(D\) 模。则
\[
 \operatorname{Hom}_D(D/DP,V)\longrightarrow\{v\in V:Pv=0\},
 \qquad f\longmapsto f(1+DP)
\]
是 \(k\) 线性同构。
\end{proposition}
\begin{proof}
左模同态由 \(1+DP\) 的像决定，而 \(P(1+DP)=0\)，故其像必须被 \(P\) 零化。反过来，对满足 \(Pv=0\) 的 \(v\)，令
\(f(Q+DP)=Qv\)。若两个代表元相差 \(RP\)，其作用差为 \(R(Pv)=0\)，所以此式良定义且为左模同态。两个构造互逆。
\end{proof}

对方程组 \(\sum_jP_{ij}y_j=0\)，取自由左模 \(D^r\) 的基 \(e_1,\ldots,e_r\)，再模去由关系 \(\sum_jP_{ij}e_j\) 生成的左子模。同样的论证把该商模到 \(V\) 的同态，识别为 \(V^r\) 中的解。因而系数作用在左边，关系也必须按左模生成，不能随意交换微分算子的次序。

\ProseParagraphBreak
例如 \(D/D\partial\) 的多项式解是常数，\(D/D(\partial-1)\) 的多项式解为零；第二个模本身却非零，在解析函数中还有 \(e^x\) 这样的解。方程模是否为零与某一函数空间里是否有非零解，是不同问题。

\ProseParagraphBreak
点支撑方程 \(xy=0\) 更能说明解对象的作用。取 \(D=A_1(\mathbb C)\)，在多项式或原点附近的全纯函数中，\(xf=0\) 都迫使 \(f=0\)。但在实直线的复值分布空间 \(\mathcal D'(\mathbb R)\) 中，乘法和求导分别定义为
\[
 (xT)(\varphi)=T(x\varphi),\qquad
 (\partial T)(\varphi)=-T(\varphi'),
\]
其中 \(\varphi\) 为光滑紧支撑测试函数；这两个作用满足 \([\partial,x]=1\)，所以分布空间也是左 \(D\) 模。原点的 Dirac 分布 \(\delta_0(\varphi)=\varphi(0)\) 非零，却满足 \(x\delta_0=0\)。\cref{b5:prop:D-solution-module}因此给出非零同态
\[
 D/Dx\longrightarrow\mathcal D'(\mathbb R),\qquad
 1+Dx\longmapsto\delta_0.
\]
Weyl 关系还给出 \(x\partial\delta_0=-\delta_0\) 及
\(x\partial^j\delta_0=-j\partial^{j-1}\delta_0\)。这些点支撑分布及其导数保存了普通函数解无法检测的信息；方程模先保存算子关系，再由所选的解对象检测它们。

\subsection{微分算子、滤过与伴随分次}
\label{b5:subsec:D01:02}

\begin{definition}\label{b5:def:D-operator-order}
设 \(k\) 为特征零域，\(R\) 为交换含幺 \(k\) 代数。在 \(\operatorname{End}_k(R)\) 中递归规定 \(F_{-1}=0\)，并令
\[
 F_m=\bigl\{P:[P,a]\in F_{m-1}\;\text{对所有}\;a\in R\bigr\}\qquad(m\ge0),
\]
其中 \(P\in\operatorname{End}_k(R)\)，\(a\) 表示乘以 \(a\) 的算子，\([P,a]=Pa-aP\)。\(\mathcal D(R)=\bigcup_{m\ge0}F_m\) 的元素称为 \(R\) 上的
\term[daishu-weifen-suanzi]{代数微分算子}[algebraic differential operator]，最小的相应 \(m\) 称为其阶。
\end{definition}
\symidx{\(\mathcal D(R)\)：代数上的微分算子环}

\(F_0=\operatorname{End}_R(R)=R\)，所以零阶算子就是乘法。恒等式
\([PQ,a]=P[Q,a]+[P,a]Q\) 归纳给出 \(F_mF_n\subset F_{m+n}\)。在 \(R=k[x]\) 中，归纳还给出
\[
 F_m=\sum_{j=0}^m k[x]\partial^j.
\]
事实上，若 \([P,x]=\sum_{j<m}b_j(x)\partial^j\)，令
\(Q=\sum_{j<m}b_j(x)\partial^{j+1}/(j+1)\)，则
\([P-Q,x]=0\)。因此 \(P-Q\) 与一切多项式乘法交换，必为乘法算子。这证明 \(\mathcal D\bigl(k[x]\bigr)=A_1(k)\)。

\ProseParagraphBreak
这个阶滤过把 \(x\) 放在零阶，把 \(\partial\) 放在一阶。其伴随分次环为
\[
 \operatorname{gr}_F D\simeq k[x,\xi],
\]
其中 \(\xi\) 是 \(\partial\) 的一阶符号；对
\(P=a_m(x)\partial^m+\cdots\)、\(a_m\ne0\)，主符号是 \(a_m(x)\xi^m\)。因为交换子降低阶，原来的非交换关系在伴随分次环中变为交换关系。

\ProseParagraphBreak
还有常用于增长估计的 Bernstein 滤过，把 \(x,\partial\) 都放在一次。它同样产生多项式分次环，但滤过不相同。例如 \(\partial-x\) 的阶主符号为 \(\xi\)，Bernstein 主符号却为 \(\xi-x\)。以下几何特征簇始终用阶滤过；不能因两个分次环抽象同构就把两个主符号混用。两种滤过的区别亦见 \cite[Section 3.2, pp. 88--90]{HottaTakeuchiTanisaki2008DModules}。

\subsection{光滑代数簇上的 D-模与特征簇}
\label{b5:subsec:D01:03}

设 \(X\) 为光滑复代数簇，\(\mathcal O_X\) 为正则函数层，\(\mathcal T_X=\operatorname{Der}_{\mathbb C}(\mathcal O_X)\) 为切向量层。层 \(\mathcal End_{\mathbb C}(\mathcal O_X)\) 的局部截面是与开集限制相容的 \(\mathbb C\) 线性层态射。把函数视为乘法算子，由 \(\mathcal O_X\) 与 \(\mathcal T_X\) 在其中生成的子代数层记为 \(\mathcal D_X\)，称为微分算子层。这一定义直接在层态射中进行，因而其算子已经具有限制映射。

\ProseParagraphBreak
局部可以选择正则函数 \(x_1,\ldots,x_n\)，使它们给出到 \(\mathbb A^n\) 的平展态射；这样的函数称为平展坐标。缩小开集后，等价的条件是 \(\mathrm{d}x_1,\ldots,\mathrm{d}x_n\) 构成余切层 \(\Omega_X^1\) 的基。其对偶导子 \(\partial_1,\ldots,\partial_n\) 满足
\[
 [\partial_i,a]=\frac{\partial a}{\partial x_i},\qquad
 [\partial_i,\partial_j]=0.
\]
把函数移到左边、把相互交换的导子按次序排列，便得到局部表达
\[
 \mathcal D_U=\bigoplus_{\alpha\in\mathbb N^n}\mathcal O_U\partial^\alpha,
 \qquad
 F_m\mathcal D_U=\bigoplus_{|\alpha|\le m}\mathcal O_U\partial^\alpha.
\]
其中直和表达的唯一性、阶滤过与递归交换子定义的一致性，以及
\(\operatorname{gr}\mathcal D_X\simeq\operatorname{Sym}_{\mathcal O_X}\mathcal T_X\)，可由以下符号计算得到。右侧是余切丛 \(T^*X\) 上沿纤维的多项式函数。
这个对偶方向来自求值：切向量场 \(v\) 给出余切丛上的函数 \((x,\xi)\mapsto\xi(v_x)\)，它在每个余切纤维上为线性函数。若 \(\xi=\sum_i\xi_i\cdot\mathrm{d}x_i\)，则 \(\xi(\partial_i)=\xi_i\)；因此导子 \(\partial_i\) 的符号恰是余切坐标 \(\xi_i\)，其对称乘积给出纤维上的多项式函数。
\begin{proofsketch}
对于由函数和导子组成的表达，反复与各个 \(x_i\) 取交换子：若最高总次数为 \(m\)，任取 \(|\alpha|=m\)，与 \(x_i\) 分别取 \(\alpha_i\) 次交换子所得的算子就是乘以 \(\alpha!a_\alpha\)。因此零算子的最高系数全为零，归纳证明正规表达唯一。交换子 \([P,a]\) 降低至少一阶，故上面的阶滤过满足\cref{b5:def:D-operator-order}的递归条件。

反过来，设层态射 \(P\) 按递归交换子条件的阶数至多为 \(m\)。其 \(m\) 重交换子是乘法算子；所乘函数关于这 \(m\) 个函数参数是对称的多导子。这由乘法算子彼此交换及 Leibniz 恒等式得到。由于 \(\mathrm{d}x_i\) 为余切层的基，该多导子由在坐标参数上的取值决定。用这些取值除以相应的 \(\alpha!\)，得到系数 \(a_\alpha\)，使
\(P-\sum_{|\alpha|=m}a_\alpha\partial^\alpha\) 的全部 \(m\) 重交换子为零，即阶数至多为 \(m-1\)。归纳后 \(P\) 属于函数与导子生成的算子层，证明两种定义一致。

令 \(\xi_i\) 为 \(\partial_i\) 的一阶符号，局部便有 \(\operatorname{gr}\mathcal D_U=\mathcal O_U[\xi_1,\ldots,\xi_n]\)。换坐标时，一阶符号按切向量变换，各阶乘积按其对称幂变换，故这些识别黏合成所示对称代数同构。局部坐标表达、阶滤过性质、递归描述和主符号识别分别见\cite[Section 1.1, Proposition 1.1.3, Exercise 1.1.4 and the principal-symbol discussion, pp. 15--17]{HottaTakeuchiTanisaki2008DModules}。
\end{proofsketch}

求出某个函数空间中的全部解，往往比记录方程的最高阶约束困难。阶滤过使最高阶关系落在交换的符号代数中，因而可以用其支撑描述约束所在的余切方向。对一个模需要同时记录全部关系，并保证所选滤过由有限数据控制；良滤过提供了这个条件。

\ProseParagraphBreak
在仿射坐标开集上，\(\operatorname{gr}\mathcal D_X\) 的截面环为 Noether 多项式环。左理想或右理想的分次理想都有有限个齐次生成元，提升这些生成元并按阶归纳，便得到原理想的有限生成元。因此局部算子环及各茎均为左、右 Noether 环。对于底层 \(\mathcal O_X\) 模拟凝聚、且在 \(\mathcal D_X\) 上局部有限生成的模，局部展示的核仍在 \(\mathcal O_X\) 上拟凝聚；在仿射开集取截面后，Noether 性给出核的有限生成元。因而 \(\mathcal D_X\) 是凝聚环层，这类模也局部有限展示。反过来，局部有限展示的 \(\mathcal D_X\) 模是有限自由 \(\mathcal D_X\) 模之间态射的余核，所以其底层 \(\mathcal O_X\) 模拟凝聚。这说明下面的有限展示定义与通常的凝聚 D-模概念一致；参见\cite[Propositions 1.4.6 and 1.4.9, pp. 26--27]{HottaTakeuchiTanisaki2008DModules}。

\begin{definition}\label{b5:def:D-good-characteristic}
设 \(X\) 为光滑复代数簇，\(\mathcal O_X\) 为正则函数层，\(\mathcal D_X\) 为 \(\mathcal O_X\) 上的代数微分算子层，\(F_i\mathcal D_X\) 为阶数不超过 \(i\) 的算子层（\(i<0\) 时取零）。\(\mathcal D_X\) 的左模层称为左 D-模；若它局部有限展示，则称为凝聚 D-模（亦称相干 D-模）。\termidx[ningju-Dmo]{凝聚 D-模}\termidx[xianggan-Dmo]{相干 D-模}设 \(M\) 为凝聚左 D-模，\((F_jM)_{j\in\mathbb Z}\) 为递增、穷尽、局部下有界的滤过。若各 \(F_jM\) 为凝聚 \(\mathcal O_X\) 模，
\(F_i\mathcal D_XF_jM\subset F_{i+j}M\)，且
\(\operatorname{gr}_F M\) 在 \(\operatorname{gr}\mathcal D_X\) 上局部有限生成，则称为
\term[liang-lvguo]{良滤过}[good filtration]。
通过 \(\operatorname{gr}\mathcal D_X\simeq\operatorname{Sym}_{\mathcal O_X}\mathcal T_X\)，其中 \(\mathcal T_X\) 为切向量层，将 \(\operatorname{gr}_F M\) 视为余切丛 \(T^*X\) 上的凝聚层。其支撑集称为 \(M\) 的
\term[tezheng-cu]{特征簇}[characteristic variety]，记为 \(\operatorname{Ch}(M)\)。
\end{definition}
\symidx{\(T^*X,\operatorname{Ch}(M)\)：余切丛与 D-模的特征簇}

\begin{proposition}[良滤过与特征簇]\label{b5:prop:D-good-filtration-invariance}
设 \(X\) 为光滑复代数簇，\(\mathcal D_X\) 为代数微分算子层，\(M\) 为凝聚左 \(\mathcal D_X\) 模。它局部具有关于阶滤过的良滤过，以其伴随分次模在 \(T^*X\) 上的支撑定义的特征簇 \(\operatorname{Ch}(M)\) 不依赖良滤过的选择。对凝聚左 \(\mathcal D_X\) 模的短正合列
\(0\to M'\to M\to M''\to0\)，有
\(\operatorname{Ch}(M)=\operatorname{Ch}(M')\cup\operatorname{Ch}(M'')\)。若 \(\pi:T^*X\to X\) 为丛投影，则还有
\[
 \operatorname{Supp}_{\mathcal O_X}(M)=\pi\bigl(\operatorname{Ch}(M)\bigr),
\]
且这个支撑是闭集。
\end{proposition}
\begin{proof}
局部取有限个 \(\mathcal D_X\) 生成元 \(u_i\)，指定它们的整数滤过次数 \(d_i\)，令
\(F_jM=\sum_iF_{j-d_i}\mathcal D_Xu_i\)。每项在 \(\mathcal O_X\) 上凝聚，分次模由这些元的符号生成，故为良滤过。

两个良滤过 \(F,G\) 的有限个分次生成元提升到模中，分别在另一滤过的有限次数出现。因此可取 \(c\ge0\)，使
\(F_jM\subseteq G_{j+c}M\)、\(G_jM\subseteq F_{j+c}M\) 对所有 \(j\) 成立。设次数 \(m\) 的齐次符号 \(a\) 零化 \(\operatorname{gr}_FM\)，取它的算子提升 \(P\)，则
\(PF_jM\subseteq F_{j+m-1}M\)。迭代 \(N\) 次给出 \(P^NF_iM\subseteq F_{i+Nm-N}M\)。对整数 \(N>2c\)，从 \(G_jM\subseteq F_{j+c}M\) 出发，得到
\[
\begin{aligned}
 P^NG_jM&\subseteq P^NF_{j+c}M
          \subseteq F_{j+c+Nm-N}M\\
        &\subseteq G_{j+Nm+2c-N}M
          \subseteq G_{j+Nm-1}M.
\end{aligned}
\]
最后一个包含用的是 \(2c-N\le-1\)。所以 \(a^N\) 零化 \(\operatorname{gr}_GM\)。交换 \(F,G\) 后得到反向包含，故两分次模的零化理想有相同根理想，支撑相同。

在 \(M'\) 上取 \(M\) 的交滤过，在 \(M''\) 上取商滤过，得到分次短正合列。分次环局部为 Noether 多项式环，所以分次子模和商模都有限生成，各滤过项也凝聚；这些仍为良滤过。交换环上短正合列的支撑是两端支撑的并，故得到结论。参见\cite[Theorems 2.1.3 and 2.2.1, pp. 58--59]{HottaTakeuchiTanisaki2008DModules}。

最后证明支撑与投影的关系。局部记 \(G=\operatorname{gr}_FM\)。滤过穷尽且下有界，故 \(M_x=0\) 当且仅当 \(G_x=0\)。\(G_x\) 在 \(\mathcal O_{X,x}[\xi_1,\ldots,\xi_n]\) 上分次有限生成，各次分量在 \(\mathcal O_{X,x}\) 上有限生成；分次 Nakayama 论证于是给出
\[
 G_x\ne0\quad\Longleftrightarrow\quad
 G_x/(\mathfrak m_x,\xi_1,\ldots,\xi_n)G_x\ne0.
\]
具体说，若右侧为零，最低非零次数的分量便等于其极大理想倍，由普通 Nakayama 引理只能为零，矛盾。右侧非零等价于 \((x,0)\in\operatorname{Ch}(M)\)。特征簇由齐次理想定义，沿纤维伸缩不变；闭性使 \((x,\xi)\) 在其中时 \((x,0)\) 也在其中。因此 \(\pi(\operatorname{Ch}(M))=\{x:(x,0)\in\operatorname{Ch}(M)\}\)，并与 \(\operatorname{Supp}_{\mathcal O_X}(M)\) 一致；这个集合是闭集。
\end{proof}

这里凝聚性是对 \(\mathcal D_X\) 而言，不要求 \(M\) 在 \(\mathcal O_X\) 上有限生成。

\begin{proposition}\label{b5:prop:D-principal-characteristic}
设 \(X=\mathbb A^1_{\mathbb C}\)，\(D=A_1(\mathbb C)\) 取阶滤过，\(\xi\) 为 \(\partial\) 的一阶符号。设非零算子 \(P=a_m(x)\partial^m+\cdots\in D\)，其中 \(m\ge0\)、\(a_m\in\mathbb C[x]\setminus\{0\}\)。对左商模 \(M=D/DP\) 的商滤过，有
\[
 \operatorname{gr}M\simeq\mathbb C[x,\xi]/\bigl(a_m(x)\xi^m\bigr),
 \qquad
 \operatorname{Ch}(M)=V\bigl(a_m(x)\xi^m\bigr).
\]
\end{proposition}
\begin{proof}
主符号的乘积满足 \(\sigma(QP)=\sigma(Q)\sigma(P)\)：最高阶系数的乘积非零，故阶相加。左理想 \(DP\) 的每个元素就是某个 \(QP\)，所以其分次理想恰由 \(\sigma(P)=a_m(x)\xi^m\) 生成。商滤过的分次模于是为所写商环，它的支撑为该理想的零点集。
\end{proof}

例如 \(D/D\partial\) 的特征簇是零截面 \(\xi=0\)，\(D/Dx\) 的特征簇是点 \(x=0\) 上的整条余切纤维。后一模非零：在 \(\mathbb C[z]\) 上令 \(\partial\) 乘以 \(z\)、\(x=-\mathrm{d}/\mathrm{d}z\)，可验证 \(\partial x-x\partial=1\)，且向量 \(1\) 被 \(x\) 零化。这给出了该商模的具体模型。

\ProseParagraphBreak
再取 \(P=x\partial-\lambda\)、\(\lambda\in\mathbb C\)。它的阶主符号为 \(x\xi\)，所以
\[
\operatorname{Ch}\bigl(D/D(x\partial-\lambda)\bigr)
=V(x\xi)=V(\xi)\cup V(x).
\]
在 \(x\ne0\) 处只留下零截面；在原点，最高阶系数消失，整条余切纤维出现。不同的 \(\lambda\) 都给出同一特征簇，而方程 \(xy'=\lambda y\) 在避开原点的单连通解析区域内有解 \(y=cx^\lambda\)，其绕圈因子随 \(\lambda\) 改变。特征簇保留了最高阶约束的几何形状，低阶参数及解析延拓还需要其他资料。

\ProseParagraphBreak
即使绕圈因子相同，原点处的延拓仍可能不同。令
\(M_0=D/D(x\partial)\)、\(M_{-1}=D/D(x\partial+1)\)，两者的特征簇都是 \(V(x\xi)\)。在穿孔圆盘上，解分别为常数与 \(c/x\)，单值均为 \(1\)，因而给出同构的局部系统。若用原点的全纯函数芽环 \(\mathcal O=\mathbb C\{x\}\) 检测它们，则前者的解空间是一维，后者只有零解：写 \(f=\sum_{j\ge0}a_jx^j\)，方程 \(xf'+f=0\) 给出 \((j+1)a_j=0\)。由\cref{b5:prop:D-solution-module}，
\[
 \dim_{\mathbb C}\operatorname{Hom}_D(M_0,\mathcal O)=1,
 \qquad
 \operatorname{Hom}_D(M_{-1},\mathcal O)=0.
\]
所以这两个代数 D-模不同构。奇点外的局部系统只描述穿孔区域中的延拓，不能决定模在奇点处允许怎样的截面。

\subsection{holonomic D-模与术语辨析}
\label{b5:subsec:D01:04}

特征簇越小，方程的约束通常越强。不过约束不能任意增加而仍保留非零凝聚模。

\begin{proposition}\label{b5:prop:D-simultaneous-annihilation}
设 \(D=A_1(\mathbb C)\)。若左 \(D\) 模中的向量 \(v\) 同时满足 \(xv=0\)、\(\partial v=0\)，则 \(v=0\)。等价地，\(D\partial+Dx=D\)。
\end{proposition}
\begin{proof}
由 Weyl 关系，\(v=(\partial x-x\partial)v=0\)。在左理想中，\(\partial x\in Dx\)、\(x\partial\in D\partial\)，所以它们的差 \(1\) 属于 \(D\partial+Dx\)，该左理想就是 \(D\)。
\end{proof}

在一维余切空间中，\(D/D\partial\) 与 \(D/Dx\) 的特征簇分别是直线 \(V(\xi)\) 与 \(V(x)\)。若只看两条生成关系的主符号，同时施加它们似乎会留下零维交点 \(x=\xi=0\)；上述命题却说明商模已为零，其特征簇为空。整个分次关系理想还包含 \(1\) 的符号，不能只看所写生成元的符号。这是\cite[Exercise 2.2.6, pp. 61--62]{HottaTakeuchiTanisaki2008DModules}所述现象的一个具体计算，也说明位置和微分方向不能在非零模中同时完全零化。一般情形中的维数下界由下面的不等式给出。

\begin{theorem}[Bernstein 不等式]\label{b5:thm:D-bernstein}
设 \(n\ge0\)，\(X\) 为纯 \(n\) 维光滑复代数簇，\(\mathcal D_X\) 为代数微分算子层，\(M\ne0\) 为凝聚左 \(\mathcal D_X\) 模，\(\operatorname{Ch}(M)\subseteq T^*X\) 为其关于阶滤过的特征簇。则
\(\dim\operatorname{Ch}(M)\ge n\)。
\end{theorem}

\begin{proofsketch}
按 \(n=\dim X\) 归纳，\(n=0\) 时结论直接成立。光滑簇的不可约分量彼此不交，可以先限制到 \(M\) 非零的一个分量并仍记它为 \(X\)。由\cref{b5:prop:D-good-filtration-invariance}，模的支撑等于特征簇的投影。若支撑为 \(X\)，特征簇便包含零截面，维数至少为 \(n\)。否则，支撑为非空真闭子集。取其中一个不可约分量的光滑点附近的开集，并去掉其他支撑分量，便可选择光滑超曲面 \(S=\{x_1=0\}\) 包含所剩支撑；限制后的 \(M\) 仍非零。取与 \(S\) 相容的平展坐标及法向算子 \(\partial_1\)。底层 \(\mathcal O_X\) 模拟凝聚，所以支撑在 \(S\) 上意味着局部化 \(M[x_1^{-1}]=0\)，每个局部元素 \(m\) 都被某个 \(x_1\) 的幂零化。

令 \(N=\ker(x_1:M\to M)\)，把它视为 \(S\) 上的层。函数在 \(N\) 上通过 \(\mathcal O_S=\mathcal O_X/(x_1)\) 作用，切向导子 \(\partial_2,\ldots,\partial_n\) 保持这个核，因而赋予它左 \(\mathcal D_S\) 模结构。对非零局部元素 \(m\)，取最小的 \(r\ge1\) 使 \(x_1^rm=0\)；则 \(x_1^{r-1}m\) 是 \(N\) 的非零元素。

所需投影可以先在 \(D/Dx\) 上理解。记 \(\delta=1+Dx\)，正规表达给出 \(D/Dx=\bigoplus_{r\ge0}\mathbb C\partial^r\delta\)，且 \(x\partial^r\delta=-r\partial^{r-1}\delta\)。故 \(\ker x=\mathbb C\delta\)，我们希望保留 \(a_0\delta\)，消去各项 \(a_r\partial^r\delta\)（\(r>0\)）。例如在 \(\partial\delta\) 上，\(\partial\delta+\partial(x\partial\delta)=0\)；更高次项的抵消由同样的 Weyl 关系控制。一般模中，\(N\) 承担这个常数项的角色，Weyl 关系给出逐元素有限求和的投影
\[
 p(m)=\sum_{j\ge0}\frac{\partial_1^j x_1^jm}{j!}\in N,
 \qquad
 m=\sum_{j\ge0}\frac{(-1)^j}{j!}\partial_1^j p(x_1^jm).
\]
这个投影是在取法向导数变量的常数项。确实，对 \(n\in N=\ker x_1\)，Weyl 关系递推给出
\(x_1^j\partial_1^r n=(-1)^j r!\,\partial_1^{r-j}n/(r-j)!\)（\(0\le j\le r\)），\(j>r\) 时为零。所以
\[
p(\partial_1^r n)
=\left(\sum_{j=0}^r(-1)^j\binom rj\right)\partial_1^r n
=\begin{cases}n,&r=0,\\0,&r>0.\end{cases}
\]
它保留 \(N\) 中的常数项，并消去全部正次法向导数。对一般 \(m\)，用 \(x_1\partial_1^j=\partial_1^jx_1-j\partial_1^{j-1}\)，可将第一式的 \(x_1\) 倍写成
\[
 x_1p(m)=\sum_{j\ge0}\frac{\partial_1^jx_1^{j+1}m}{j!}
 -\sum_{j\ge1}\frac{\partial_1^{j-1}x_1^jm}{(j-1)!}=0.
\]
两和重编号后逐项相消。将 \(p(x_1^jm)\) 的定义代入第二式，\(\partial_1^rx_1^rm\) 的系数为
\[
\sum_{j=0}^r\frac{(-1)^j}{j!(r-j)!}
=\frac{1}{r!}\sum_{j=0}^r(-1)^j\binom rj.
\]
由二项式公式，它在 \(r>0\) 时为零，在 \(r=0\) 时为一，所以其和确为 \(m\)。若 \(\sum_{j=0}^s\partial_1^jn_j=0\)、\(n_j\in N\)，作用 \(x_1^s\) 得 \((-1)^ss!n_s=0\)，再按 \(s\) 下降归纳。因此这些和唯一，局部有
\[
 M\simeq\mathbb C[\partial_1]\otimes_{\mathbb C}i_*N,
 \qquad i:S\hookrightarrow X,
\]
其中 \(x_1\) 按 \(-\mathrm{d}/\mathrm{d}\partial_1\) 作用。这是 Kashiwara 等价在超曲面支撑情形的局部模型，参见\cite[Theorem 1.6.1 and its proof, pp. 48--50]{HottaTakeuchiTanisaki2008DModules}。

为核对凝聚性，取有限个局部 \(\mathcal D_X\) 生成元 \(m_1,\ldots,m_q\)，并选择 \(r_i\ge1\) 使 \(x_1^{r_i}m_i=0\)。令 \(N_0\) 为有限个元素
\(p(x_1^jm_i)\)、\(1\le i\le q\)、\(0\le j<r_i\)，在 \(\mathcal D_S\) 上生成的子模，并置
\(W=\sum_{j\ge0}\partial_1^ji_*N_0\)。Taylor 式使所有 \(m_i\) 属于 \(W\)。这个子模对 \(x_1\)、\(\partial_1\) 及切向导子封闭；对任意局部函数 \(a\)，恒等式
\[
 a\partial_1^jn
 =\sum_{\ell=0}^j(-1)^\ell\binom{j}{\ell}
   \partial_1^{j-\ell}\bigl((\partial_1^\ell a)n\bigr)
 \qquad(n\in N_0)
\]
也证明它对函数作用封闭，其中括号内的函数在 \(S\) 上限制后作用于 \(N_0\)。因此 \(W=M\)。上述唯一分解又给出 \(W\cap\ker x_1=i_*N_0\)，故 \(N=N_0\) 在 \(\mathcal D_S\) 上局部有限生成。其底层 \(\mathcal O_S\) 模拟凝聚，因为它是 \(\mathcal O_X\) 拟凝聚模之间映射的核且被 \(x_1\) 零化；由前面的 Noether 性与凝聚性结论，\(N\) 是凝聚左 \(\mathcal D_S\) 模。

给 \(N\) 取良滤过 \(G\)，在 \(M\) 上令
\[
 F_jM=\sum_{r\ge0}\partial_1^ri_*G_{j-r}N.
\]
每个固定 \(j\) 只有有限个非零项。上面的函数作用公式与法向关系说明这是良滤过，并有
\[
 \operatorname{gr}_FM\simeq
 \mathbb C[\xi_1]\otimes_{\mathbb C}i_*\operatorname{gr}_GN.
\]
这里 \(x_1\) 的符号作用为零，\(\xi_1\) 是自由的法向余切坐标；切向符号按 \(\operatorname{gr}_GN\) 的原作用给出。因此
\(\dim\operatorname{Ch}(M)=\dim\operatorname{Ch}(N)+1\)。归纳假设给出 \(\dim\operatorname{Ch}(N)\ge n-1\)，于是 \(\dim\operatorname{Ch}(M)\ge(n-1)+1=n\)，归纳完成。局部模型在重叠坐标上的层论识别由 Kashiwara 等价给出；闭嵌入的特征簇公式及上述支撑归纳见\cite[Lemma 2.3.5 and the proof of Corollary 2.3.2, pp. 63--64]{HottaTakeuchiTanisaki2008DModules}。
\end{proofsketch}

设 \(X\) 为纯 \(n\) 维光滑复代数簇，\(\mathcal D_X\) 为代数微分算子层。零左 \(\mathcal D_X\) 模，以及满足 \(\dim\operatorname{Ch}(M)=n=\dim X\) 的非零凝聚左 \(\mathcal D_X\) 模，称为
\term[holonomic-D-mo]{holonomic D-模}[holonomic D-module]。
Bernstein 不等式说明，这正是非零凝聚 D-模能够达到的最小特征簇维数。因此\cref{b5:prop:D-principal-characteristic}中，凡商模非零且 \(P\ne0\)，其特征簇是平面中维数一的闭集，可由若干不可约曲线组成，所得模便是 holonomic；反之自由模 \(D\) 的特征簇是整个二维余切丛，故不是 holonomic。

\ProseParagraphBreak
在二维空间中，单个微分方程已经未必有这种最小维数。取 \(X=\mathbb A^2_{\mathbb C}\)，其 Weyl 代数记为
\(D_2=\mathbb C[x,y]\langle\partial_x,\partial_y\rangle\)，余切坐标为 \(\xi_x,\xi_y\)。正规表达表明，左理想 \(D_2\partial_x\) 由微分单项式含正次 \(\partial_x\) 的项组成，所以
\[
 \operatorname{gr}(D_2/D_2\partial_x)
 =\mathbb C[x,y,\xi_x,\xi_y]/(\xi_x).
\]
这个特征簇有三维，因而商模不是 holonomic；其局部全纯解为任意 \(g(y)\)。再添关系 \(\partial_y\) 后，商模同构于 \(\mathbb C[x,y]\)：函数作用是乘法，两偏导按通常求导作用；正规表达把其余各项都消去。此时特征簇为二维零截面 \(V(\xi_x,\xi_y)\)，模为 holonomic，局部解只剩常数。新增的独立微分约束确实缩小了允许的余切方向；但 holonomic 的定义仍是特征簇的维数条件，不能在奇点处直接用普通函数解空间的维数替代。

\ProseParagraphBreak
有限秩向量丛上的平坦联络给出另一类例子。在局部基中，把 \(\partial_i\) 作用写成 \(\nabla_i=\partial_i+B_i\)，Leibniz 法则处理它与函数的关系，而
\[
 [\nabla_i,\nabla_j]=0
 \quad\Longleftrightarrow\quad
 \partial_iB_j-\partial_jB_i+[B_i,B_j]=0
\]
保证各偏导作用满足 D-模关系。这就是平坦性。取 \(F_jM=M\) 对 \(j\ge0\)、\(F_jM=0\) 对 \(j<0\)，便得良滤过，其正次符号全部作用为零，特征簇为零截面，所以非零联络是 holonomic。一般 holonomic 模却未必是向量丛上的联络，\(D/Dx\) 就说明了点支撑的可能性。

\ProseParagraphBreak
英文 holonomic 表示这里的特征簇维数条件，holomorphic 则表示复解析意义的“全纯”，二者不能互译。D-模与第三卷的微分代数也有不同的研究对象：前者研究线性微分算子作用及其模，后者还允许非线性微分多项式关系；两者的联系须在具体对象上讨论。

% !TeX root = ../../Book5.tex
% 纲要标题：### D.2 代数、形式与解析解
% 纲要位置：工作记录/计划纲要.md，第 5041 行。
% 纲要细目（保留原定写作范围）：
%
% * 有理函数域上的基本解矩阵、无新常数与微分自同构群就地说明；PV 存在性与 Galois 对应在附录 E 建立
% * 形式局部方程的系数递推、Puiseux 扩张与解析单值；整数指数平移改变对数格，Laurent 两项复形与圆周上同调直接比较
% * 正则奇点、非正则奇点和 Stokes 数据所需的额外分析准备；区分积分平面与解变量扇形，齐次化算例给出明确的 Stokes 矩阵
% * Riemann–Hilbert 对应的适用对象与条件，并给出联络及导出版本的证明概要；用 jet 线性分裂解释 GAGA 输入，用常数、曲线及点支撑模说明 de Rham 移位，并核对无穷远正则性
% * 与分析学、代数几何和几何表示论的联系，不把全部解析理论压入本卷
% #### D.2.1 有理函数域与形式局部方程
% #### D.2.2 奇点与 Stokes 数据的分析前提
% #### D.2.3 Riemann–Hilbert 对应简介
% #### D.2.4 与分析学、几何表示论的联系
% ---
%
\section{代数、形式与解析解}
\label{b5:sec:D02}

同一个有理系数方程可以在函数域、形式 Laurent 级数域，或复解析函数空间中求解。前两者关心相应域中的代数关系，第三种解释还记录沿路径延拓和接近奇点时的增长。以下以一阶方程说明这些信息为何不能相互替代。

\subsection{有理函数域与形式局部方程}
\label{b5:subsec:D02:01}

在 \(K=\mathbb C(x)\) 上取通常求导，设 \(A(x)\in M_n(K)\)、\(n\ge1\)。若把求导延拓到扩域 \(L\)，其中满足 \(Z'=A(x)Z\) 的可逆矩阵 \(Z\in\operatorname{GL}_n(L)\) 称为基本解矩阵。若 \(L=K(Z_{ij})\)，且 \(L\) 中导数为零的元素仍只有 \(\mathbb C\)，则称 \(L/K\) 为这个系统的 Picard--Vessiot 扩张。固定 \(K\) 并与求导交换的 \(L\) 的域自同构组成其微分 Galois 群\footnote{伽罗瓦（Évariste Galois，1811-1832），法国数学家，以方程根的置换研究根式可解性，奠定了群论与域论的重要基础。}；\appref{b5:app:E}建立相应的存在性与 Galois 对应。

\ProseParagraphBreak
若选定有限点 \(a\)，令 \(z=x-a\)，有理函数又可以展开到 \(\mathbb C(\!(z)\!)\) 中。同一个系统因此也可以在形式 Laurent 级数域上研究；基域变大以后，解之间的代数关系和相应的微分 Galois 群一般都会改变。在常点，基本解矩阵可以直接由系数递推得到。

\begin{proposition}\label{b5:prop:D-ordinary-formal-solution}
设 \(n\ge1\)，\(A(z)=\sum_{j\ge0}A_jz^j\in M_n\bigl(\mathbb C[\![z]\!]\bigr)\)，\(Z_0\in\operatorname{GL}_n(\mathbb C)\)，撇号表示形式求导 \(\mathrm d/\mathrm dz\)。则存在唯一
\(Z(z)\in\operatorname{GL}_n\bigl(\mathbb C[\![z]\!]\bigr)\) 满足 \(Z'=AZ\) 且 \(Z(0)=Z_0\)。
\end{proposition}
\begin{proof}
写 \(Z=\sum_{m\ge0}Z_mz^m\)，比较系数得到递推
\[
 (m+1)Z_{m+1}=\sum_{j=0}^m A_jZ_{m-j}.
\]
特征零保证可以除以 \(m+1\)，所以给定 \(Z_0\) 后依次唯一确定各系数；反向代入保证方程成立。行列式的常数项为 \(\det Z_0\ne0\)，故 \(Z\) 在幂级数环上可逆。
\end{proof}

该命题给出形式解的存在唯一性，级数的收敛性还须另证。当 \(A\) 在圆盘内解析时，解析常微分方程的存在唯一性给出收敛解，并与上述形式解相符。穿孔圆盘中，解可以沿闭路延拓。记延拓后的基本解矩阵为 \(\widetilde Z\)，它仍满足 \(\widetilde Z'=A\widetilde Z\)。令 \(M=Z^{-1}\widetilde Z\)，则
\[
 M'=-Z^{-1}AZZ^{-1}\widetilde Z+Z^{-1}A\widetilde Z=0.
\]
因此 \(M\in\operatorname{GL}_n(\mathbb C)\)，且 \(\widetilde Z=ZM\)。这就是
\term[danzhi-juzhen]{单值矩阵}[monodromy matrix]，换基本解矩阵只改变其共轭类。

\ProseParagraphBreak
例如在形式 Laurent 域中求 \(zy'=\lambda y\) 的解，写
\(y=\sum_{m\ge N}a_mz^m\)，比较系数便得
\[
 (m-\lambda)a_m=0\qquad(m\ge N).
\]
若 \(\lambda\notin\mathbb Z\)，形式域 \(\mathbb C(\!(z)\!)\) 中只有零解；若 \(\lambda\in\mathbb Z\)，全部解为 \(cz^\lambda\)。但在穿孔圆盘的通用覆盖上，对任何复数 \(\lambda\) 都有非零解析解
\(z^\lambda=\exp(\lambda\log z)\)，一次正向绕原点后乘以 \(e^{2\pi i\lambda}\)。例如 \(\lambda=1/2\) 的解在二次形式扩张 \(\mathbb C(\!(z^{1/2})\!)\) 中出现，并有单值 \(-1\)。一般有理数指数可在某个有限 Puiseux 扩张 \(\mathbb C(\!(z^{1/r})\!)\) 中得到；\(\lambda\notin\mathbb Q\) 时则不可能，因为该域的幂指数都在 \(\frac1r\mathbb Z\) 中，仍可用上述系数比较排除。这里形式解的存在、有限代数分支和通用覆盖上的解析解是不同问题。单值描述解析延拓，还不单独决定非正则奇点的渐近行为。

\subsection{奇点与 Stokes 数据的分析前提}
\label{b5:subsec:D02:02}

单值矩阵只记录绕圈前后的变化，尚未说明解接近奇点时怎样增长。为限制这种增长，可以寻找一个局部基，使联络矩阵至多有一阶极点。基的选择会改变极点的外观，因此应把条件写成稳定格的存在性。

\begin{definition}\label{b5:def:D-regular-singular}
设 \(\mathcal O=\mathbb C\{z\}\) 为原点处收敛幂级数环，\(K=\mathcal O[z^{-1}]\) 为亚纯函数芽域，\(M\) 为有限维 \(K\) 向量空间，配有沿 \(\partial_z=\mathrm d/\mathrm dz\) 的联络 \(\nabla_{\partial_z}\)。若存在有限自由 \(\mathcal O\) 子模 \(\Lambda\subseteq M\)，使 \(K\otimes_{\mathcal O}\Lambda\simeq M\)，并满足
\[
 z\nabla_{\partial_z}(\Lambda)\subset\Lambda,
\]
则称该联络在 \(z=0\) 有\term[zhengze-qidian]{正则奇点}[regular singularity]；否则称为非正则奇点。
\end{definition}
\symidx{\(\mathbb C\{z\}\)：原点处的收敛幂级数环}
\symidx{\(\nabla_{\partial_z}\)：沿局部求导方向的联络算子}

在格的基下，条件表示联络矩阵至多有一阶极点。它允许先作亚纯换基，因此不能只根据某个任意基中出现的高阶极点下结论。此局部条件及等价刻画见
\cite[Theorem 5.1.4 and Definition 5.1.6, pp. 131--133]{HottaTakeuchiTanisaki2008DModules}。

\ProseParagraphBreak
秩一情形可以直接检验。若 \(y'=a(z)y\)，其中
\(a(z)=\lambda/z+b(z)\)、\(b\) 全纯，则解是
\(z^\lambda\exp\int b(z)\,\mathrm{d}z\)，故为正则奇点。若 \(a\) 有真正的二阶或更高阶极点，则任何非零亚纯换基 \(y=g(z)v\) 只把系数改为 \(a-g'/g\)；写 \(g=z^mh\)、\(h(0)\ne0\)，便有 \(g'/g=m/z+h'/h\)，它只能改变一阶极点和全纯部分。因此高阶极点不能消去。特别地，\(y'=-y/z^2\) 的解 \(e^{1/z}\) 单值，却在不同方向上有急剧不同的增长，是非正则情形。

\begin{proposition}\label{b5:prop:D-integer-exponent-shift}
设 \(\lambda\in\mathbb C\)、\(m\in\mathbb Z\)。在穿孔复圆盘上，亚纯换基 \(y=z^mv\) 把方程 \(y'=\lambda y/z\) 变为 \(v'=(\lambda-m)v/z\)。两方程的单值相同，而在原点选定的全纯格可以不同。
\end{proposition}
\begin{proof}
代入 \(y'=mz^{m-1}v+z^mv'\) 并除以 \(z^m\)，得到所写新方程。两单值为 \(e^{2\pi i\lambda}\) 与 \(e^{2\pi i(\lambda-m)}\)，因为 \(m\) 是整数，它们相等。以新基 \(z^m\) 选取的格是 \(z^m\mathcal O\)，其中 \(\mathcal O=\mathbb C\{z\}\)；它在亚纯域中生成同一个一维空间，但 \(m\ne0\) 时不是原格 \(\mathcal O\)。
\end{proof}

例如指数 \(7/3\) 与 \(1/3\) 由 \(y=z^2v\) 联系，给出同一个单值。为了从单值选择一个确定的对数格，可以要求方程指数的实部落在 \([0,1)\) 中，便选中 \(1/3\) 而排除其其他整数平移。这里指数 \(\lambda\) 与联络 \(\nabla=\mathrm{d}-\lambda\cdot\mathrm{d}z/z\) 的残数符号相反：残数为 \(-\lambda\)。固定指数区间，也就固定了相应的残数代表。

\ProseParagraphBreak
正则奇点也允许对数项。取
\[
Y'=\frac NzY,\qquad
N=\begin{pmatrix}0&1\\0&0\end{pmatrix},\qquad
Z=I+N\log z=\begin{pmatrix}1&\log z\\0&1\end{pmatrix}.
\]
因 \(N^2=0\)，有 \(Z'=NZ/z\) 且 \(\det Z=1\)。联络
\(\nabla_{\partial_z}=\partial_z-N/z\) 的标准格 \(\mathcal O^2\) 被 \(z\nabla_{\partial_z}\) 保持，所以原点为正则奇点。正向绕圈后，基本解矩阵变为
\(Z(I+2\pi iN)\)，单值矩阵非平凡；解中的对数增长仍符合正则条件。正则性因而既不要求解单值，也不要求解都是亚纯函数。

\begin{example}[发散形式解与扇形解]\label{b5:exm:D-euler-stokes}
方程
\[
 x^2y'-y=-x
\]
有唯一形式幂级数解
\[
 \widehat y=\sum_{m\ge1}(m-1)!\,x^m,
\]
该级数的收敛半径为零。它仍能描述不同扇形解析解的共同渐近展开，而这些解析解可能相差非零的 \(ce^{-1/x}\)。
\end{example}
\begin{proof}
比较常数项得 \(a_0=0\)，比较一次项得 \(a_1=1\)，随后递推
\(a_m=(m-1)a_{m-1}\)，得到所写级数。相邻系数比为 \(m\)，故正的收敛半径不可能存在。

先固定从 \(0\) 沿正实轴趋向无穷的两条路径：\(\gamma_+\) 从上方绕开 \(t=1\)，\(\gamma_-\) 从下方绕开它。例如可用以 \(1\) 为中心、半径 \(0<\rho<1/2\) 的上下半圆绕行；两条路径都落在某个 \(|\arg t|\le\eta<\pi/2\) 的区域内。取 \(0<\varepsilon<\pi/2-\eta\) 并固定 \(r_0>0\)，在
\[
 0<|x|<r_0,\qquad |\arg x|\le\pi/2-\eta-\varepsilon
\]
这一闭扇形中定义
\[
 y_\pm(x)=\int_{\gamma_\pm}\frac{e^{-t/x}}{1-t}\,\mathrm{d}t.
\]
积分在无穷远指数收敛，绕极点的有限路径上被积函数解析，因此在该扇形内可对参数求导。利用 \(t/(1-t)=-1+1/(1-t)\)，得
\[
 x^2y_\pm'=y_\pm-\int_{\gamma_\pm}e^{-t/x}\,\mathrm{d}t=y_\pm-x,
\]
其中最后的积分由原函数 \(-xe^{-t/x}\) 算出。

在这些固定路径上，统一有
\[
 \operatorname{Re}(t/x)
 \ge (\sin\varepsilon)\frac{|t|}{|x|}.
\]
取 \(0<\delta<1-\rho\)。路径上 \(|t|\ge\delta\) 的部分满足 \(\operatorname{Re}(t/x)\ge(\sin\varepsilon)\delta/|x|\)。绕行段又有 \(|1-t|=\rho\) 且长度有限，无穷射线上的分母也可直接控制，故这两部分的积分都为 \(O(e^{-c/|x|})\)，其中 \(c>0\) 可固定，估计对闭扇形中的 \(x\) 一致。对 \(N\ge1\) 和 \(0\le t\le\delta\)，使用
\[
 \frac1{1-t}=\sum_{j=0}^{N-1}t^j+\frac{t^N}{1-t}.
\]
余项积分的绝对值不超过
\[
 \frac1{1-\delta}\int_0^\delta t^N
   e^{-(\sin\varepsilon)t/|x|}\,\mathrm{d}t
 \le \frac{N!}{1-\delta}
       \left(\frac{|x|}{\sin\varepsilon}\right)^{N+1}.
\]
各单项式沿完整路径的 Laplace 积分由分部积分给出 \(j!x^{j+1}\)，而从 \(\delta\) 起的尾项仍指数衰减。因此两解都有 \(\widehat y\) 作为渐近展开。按上述上、下绕行约定，\(\gamma_+-\gamma_-\) 顺时针围住极点一次，被积函数在 \(t=1\) 的留数为 \(-e^{-1/x}\)，所以
\[
 y_+(x)-y_-(x)=2\pi i e^{-1/x}.
\]
这也与齐次方程 \(x^2h'-h=0\) 的解相符。
\end{proof}

\cref{b5:fig:D-stokes-planes}把这两套几何分开：\(t\) 是积分变量，路径须绕过 \(t=1\)；\(x\) 是解的自变量，扇形限制则保证积分尾部衰减与渐近估计一致。
% !TeX root = ../Book5.tex
\begin{figure}[htbp]
	\centering
	\begin{tikzpicture}[x=1cm,y=1cm,every node/.style={font=\small},line width=.6pt]
		\begin{scope}
			\node at (2.5,1.65) {积分变量的 \(t\)-平面};
			\draw[->] (0,-1.15)--(0,1.25) node[above] {\(\operatorname{Im}t\)};
			\draw[->] (-.25,0)--(5.25,0) node[right] {\(\operatorname{Re}t\)};
			\node[below left] at (0,0) {\(0\)};
			\draw[blue!70!black,->] (0,.04)--(1.8,.04);
			\draw[blue!70!black] (1.8,0) arc[start angle=180,end angle=0,radius=.5];
			\draw[blue!70!black,->] (2.8,.04)--(5.05,.04);
			\draw[red!70!black,->] (0,-.04)--(1.8,-.04);
			\draw[red!70!black] (1.8,0) arc[start angle=180,end angle=360,radius=.5];
			\draw[red!70!black,->] (2.8,-.04)--(5.05,-.04);
			\draw[blue!70!black,->] (2.3,.5)--(2.38,.493);
			\draw[red!70!black,->] (2.3,-.5)--(2.38,-.493);
			\fill (2.3,0) circle (2pt);
			\node[below] at (2.3,-.08) {\(1\)};
			\node[blue!70!black] at (2.3,.83) {\(\gamma_+\)};
			\node[red!70!black] at (2.3,-.83) {\(\gamma_-\)};
			\node at (2.5,-1.55) {\(\gamma_+-\gamma_-\) 绕极点顺时针闭合};
		\end{scope}
		\begin{scope}[xshift=7.25cm]
			\node at (1,1.65) {解变量的 \(x\)-平面};
			\fill[blue!10] (0,0)--(58:1.7) arc[start angle=58,end angle=-58,radius=1.7]--cycle;
			\draw[->] (-.6,0)--(2.7,0) node[right] {\(\operatorname{Re}x\)};
			\draw[->] (0,-1.35)--(0,1.3) node[above] {\(\operatorname{Im}x\)};
			\draw[dashed] (0,0)--(58:1.7) (0,0)--(-58:1.7);
			\draw (0.45,0) arc[start angle=0,end angle=58,radius=.45];
			\node at (35:.75) {\(\theta\)};
			\draw[blue!70!black,->] (18:1.5)--(18:.72);
			\node[below] at (1.5,.12) {\(x\to0\)};
			\node[below left] at (0,0) {\(0\)};
			\node at (1,-1.55) {\(\theta=\pi/2-\eta-\varepsilon\)};
		\end{scope}
	\end{tikzpicture}
	\caption[Stokes 算例中的两个复平面]{Stokes 算例中的两个复平面。左图的路径绕开积分极点；右图的扇形使沿这些路径的 \(e^{-t/x}\) 衰减，且 \(e^{-1/x}\) 在 \(x\to0\) 时指数衰减。}
	\label{b5:fig:D-stokes-planes}
\end{figure}


这个例子展示了 \term[Stokes-xianxiang]{Stokes 现象}[Stokes phenomenon]所涉及的信息\footnote{斯托克斯（George Gabriel Stokes，1819-1903），爱尔兰数学家、物理学家，研究流体力学、光学和渐近级数。}：同一发散形式展开可以对应不同扇形上的解析解，重叠处的差别未被任何有限阶幂级数近似记录。把非齐次方程添加一个常数坐标，就得到齐次系统
\[
 x^2Y'=\begin{pmatrix}1&-x\\0&0\end{pmatrix}Y.
\]
它在两种绕行约定下的基本解矩阵可取为
\[
 Z_\pm=\begin{pmatrix}e^{-1/x}&y_\pm\\0&1\end{pmatrix},
 \qquad
 Z_+=Z_-S,\qquad
 S=\begin{pmatrix}1&2\pi i\\0&1\end{pmatrix}.
\]
两矩阵行列式均为 \(e^{-1/x}\ne0\)，而第二列的差正是第一列的 \(2\pi i\) 倍。这个非平凡幺幂常矩阵记录了侧向解的 Stokes 跳跃。在任意闭扇形 \(|\arg x|\le\theta<\pi/2\) 中，\(|e^{-1/x}|\le e^{-\cos\theta/|x|}\)，故该修正比每个固定幂次都小；越过 \(\operatorname{Re}(1/x)=0\) 的方向，其增长行为就会改变。一般线性系统的 Stokes 矩阵还须指定形式正规形及相容的扇形渐近归一化。相应解析条件和局部 Galois 群关系参阅
\cite[Section 1.4, pp. 35--38]{Singer2009DifferentialGalois}，所引节号与页码依作者公开稿。形式微分域上的方程还需补充相应的解析条件，才能讨论这些数据。

\subsection{Riemann–Hilbert 对应简介}
\label{b5:subsec:D02:03}

在没有奇点的复流形上，平坦全纯联络的水平截面局部组成有限维常数层，即局部系统。若 \(X\) 连通并选基点，沿路径的解析延拓给出 \(\pi_1(X)\) 的有限维表示。反向则把通用覆盖上的平凡向量丛按该表示黏合，得到平坦向量丛；这说明局部系统为何是线性微分方程的自然拓扑对象。

\begin{theorem}[正则联络的 Riemann--Hilbert 对应]\label{b5:thm:D-rh-connections}
设 \(X\) 为光滑复代数簇，\(X^{\mathrm{an}}\) 为其解析化。\(X\) 上有限秩代数向量丛的正则可积联络组成的范畴，通过解析化和取水平截面，与 \(X^{\mathrm{an}}\) 上的有限秩复局部系统范畴等价。这里正则性是指：存在一个光滑完备紧化 \(X\hookrightarrow\overline X\)，使边界为简单正规交叉除子，且联络在边界附近局部存在对数联络格。该条件与紧化的选择无关，等价于在每个这样的紧化上满足同一条件。
\end{theorem}

\begin{proofsketch}
平坦解析联络与局部系统之间的对应已由节首的路径延拓和黏合构造给出；所需的新步骤是得到代数联络，以及证明解析同态也是代数的。

先设 \(X\) 为准射影簇，取光滑射影紧化 \(\overline X\)，经消解奇点使边界 \(D\) 成为简单正规交叉除子。给定局部系统，在边界附近的穿孔坐标图中，各坐标环路的单值矩阵 \(M_i\) 彼此交换。选择特征值实部落在固定长度一的半开区间中的对数 \(A_i\)，使
\(e^{2\pi i A_i}=M_i\)。由\cref{b5:prop:D-integer-exponent-shift}可见，指数改变整数不影响单值，却会改变所选格；固定区间正是为了消除这种歧义。矩阵情形还要处理 Jordan 块，这些对数可取为 \(M_i\) 的函数，故仍彼此交换。局部模型
\[
 \nabla=\mathrm{d}-\sum_i A_i\cdot\frac{\mathrm{d}z_i}{z_i},\qquad
 Z=\prod_i z_i^{A_i}
\]
具有指定单值群，沿 \(z_i=0\) 的联络残数为 \(-A_i\)，并给出至多对数极点的联络格。固定这些代表的选择使格在坐标交上相容，从而黏合成 \(\overline X^{\mathrm{an}}\) 上的对数联络 \((\widetilde E,\widetilde\nabla)\)。

这里使用的 GAGA 定理断言：射影复代数簇上的凝聚代数层与其解析凝聚层等价，且层之间的同态也由解析化唯一对应；它的证明作为外部输入。联络满足 \(\nabla(fs)=\mathrm{d}f\otimes s+f\nabla s\)，所以不是 \(\mathcal O\) 线性映射；GAGA 的同态比较处理的是凝聚层之间的 \(\mathcal O\) 线性同态，不能直接用在 \(\nabla\) 上。一阶 jet（主部）记录截面的取值及一阶变化；其序列把 \(E\) 与 \(\Omega^1\otimes E\) 连接起来。对向量丛 \(E\)，把普通一阶 jet 序列沿 \(\Omega^1\hookrightarrow\Omega^1(\log D)\) 推出，得到对数 jet 序列
\[
 0\longrightarrow\Omega^1_{\overline X}(\log D)\otimes E
 \xrightarrow{\iota}J^1_{\log D}(E)
 \longrightarrow E\longrightarrow0.
\]
其典范 \(\mathbb C\) 线性映射 \(j^1:E\to J^1_{\log D}(E)\) 满足
\(j^1(as)=a j^1(s)+\iota(\mathrm{d}a\otimes s)\)。于是对数联络 \(\nabla\) 等价于 \(\mathcal O\) 线性分裂
\[
 \sigma=j^1-\iota\nabla:E\longrightarrow J^1_{\log D}(E).
\]
局部秩一计算说明了为何两项相减。若 \(\nabla=\mathrm{d}-\omega\)，其中 \(\omega\) 是对数一形式，在平凡基给出的 jet 坐标中有 \(j^1(f)=(f,\mathrm{d}f)\)、\(\iota(\eta)=(0,\eta)\)，因而
\[
 \sigma(f)=(f,\mathrm{d}f)-(0,\mathrm{d}f-f\omega)
          =(f,f\omega).
\]
导数项抵消，故 \(\sigma(gf)=g\sigma(f)\)。一般向量丛上同样由两项的 Leibniz 公式相减得到线性分裂；jet 序列把联络编码成了 GAGA 可以处理的同态。
先由 GAGA 将 \(\widetilde E\) 代数化为 \(E\)，再将解析分裂 \(\widetilde\sigma\) 代数化为上述分裂；jet 序列与解析化相容，便得到代数对数联络 \(\nabla\)。延拓联络到一形式上，规定
\[
 \nabla^1(\omega\otimes s)
 =\mathrm{d}\omega\otimes s-\omega\wedge\nabla s.
\]
曲率 \(\nabla^1\nabla:E\to\Omega^2_{\overline X}(\log D)\otimes E\) 是 \(\mathcal O\) 线性映射，其解析化为零；GAGA 的忠实性保证它本身为零。因此代数化保留平坦性，限制到 \(X\) 得到所需正则代数联络。

对两个正则联络之间的水平解析同态，局部基本解的增长至多是有限次幂乘对数幂；其逆矩阵也满足同类估计。因此同态矩阵在边界附近只有有限阶增长。逐坐标用 Laurent 展开和 Cauchy 系数估计，有限阶增长的单值全纯函数只有有限阶极点，所以该同态亚纯延伸到 \(\overline X\)。选一个足够大的极点界后，又由 GAGA 得到代数同态。于是函子全忠实；忠实性也可从解析化或水平局部基直接看出。

对一般光滑 \(X\)，在有限仿射开覆盖上进行上述构造；交集仍准射影，其上的水平解析同构由刚证明的全忠实性代数化，余圈等式由忠实性保留，所得联络因而在 \(X\) 上黏合。正则性也可由沿光滑代数曲线拉回的判据核验；紧化条件的选择独立性见\cite[Proposition 5.3.6 and Theorem 5.3.7, pp. 154--155]{HottaTakeuchiTanisaki2008DModules}。

这里省略光滑正规交叉紧化的存在、对数格的唯一黏合及 GAGA 的证明。对数格的构造与唯一性见\cite[Theorem 5.2.17, pp. 145--147]{HottaTakeuchiTanisaki2008DModules}；代数化及联络对应见\cite[Theorem 5.3.8, Corollary 5.3.10 and Lemma 5.3.12, pp. 155--159]{HottaTakeuchiTanisaki2008DModules}，其中 Lemma 5.3.12 用一阶主部实现上述 \(\mathcal O\) 线性改写。这给出\term[Riemann-Hilbert-duiying]{Riemann--Hilbert 对应}[Riemann--Hilbert correspondence]的联络版本。
\end{proofsketch}

无穷远条件不可省略。令 \(D=A_1(\mathbb C)\)，比较 \(D/D\partial\) 与 \(D/D(\partial-1)\)。两者的阶主符号都是 \(\xi\)，特征簇都是零截面；其方程 \(y'=0\)、\(y'=y\) 在复直线上也都给出平凡秩一局部系统，解析换基为 \(y=e^xv\)。但两个模作为 \(\mathbb C[x]\) 模均自由且秩一，代数同构只能由单位 \(g\in\mathbb C[x]^\times=\mathbb C^\times\) 给出。常数换基不能改变 \(\partial\) 的作用，故它们作为代数 D-模不同构；\(e^x\) 不是代数可逆函数。在 \(z=1/x\) 坐标下，\(y'=y\) 化为 \(\mathrm{d}y/\mathrm{d}z=-y/z^2\)，由上一小节的秩一判别可知它在无穷远非正则；常数方程则仍为正则。因此相同特征簇和解析局部系统不足以分类全部代数 holonomic 模，定理中的无穷远正则性正排除了这个反例。

\ProseParagraphBreak
若允许奇点处的支撑和延拓，局部系统已不够，需要构造性层复形。
\begin{theorem}[导出 Riemann--Hilbert 对应]\label{b5:thm:D-rh-derived}
设 \(X\) 为纯 \(n\) 维光滑复代数簇，\(\mathcal D_X\) 为代数微分算子层，\(X^{\mathrm{an}}\) 为其解析化。记 \(\Omega^p=\Omega^p_{X^{\mathrm{an}}}\)，以下张量积取在 \(\mathcal O_{X^{\mathrm{an}}}\) 上。对左 \(\mathcal D_X\) 模 \(M\)，采用 de Rham 复形\footnote{德拉姆（Georges de Rham，1903-1990），瑞士数学家，研究微分拓扑与代数拓扑，建立了微分形式上同调与拓扑上同调的联系。}的约定
\[
 \operatorname{DR}_X(M)
 =\left[
 M^{\mathrm{an}}\longrightarrow
 \Omega^1\otimes M^{\mathrm{an}}
 \longrightarrow\cdots\longrightarrow
 \Omega^n\otimes M^{\mathrm{an}}
 \right][n],
\]
未移位的各项位于次数 \(0,\ldots,n\)；按上链约定 \((C[n])^j=C^{j+n}\)，移位后位于 \(-n,\ldots,0\)。局部微分由左模作用给出的联络 \(\nabla m=\sum_i\mathrm{d}z_i\otimes\partial_i m\) 决定：在 \(p\) 形式上为
\[
 \mathrm{d}_{\nabla}(\omega\otimes m)
 =\mathrm{d}\omega\otimes m+(-1)^p\omega\wedge\nabla m.
\]
对有界左模复形 \(M^\bullet\)，先对 \(\Omega^p\otimes M^{q,\mathrm{an}}\) 作总次数 \(p+q\) 的总化，总微分为 \(\mathrm{d}_{\nabla}+(-1)^p\mathrm{d}_{M}\)，再移位 \([n]\)；具体的 de Rham 复形见\cite[Section 4.2, pp. 103--105]{HottaTakeuchiTanisaki2008DModules}。这一函子给出范畴等价
\[
 \operatorname{DR}:D^b_{\mathrm{rh}}(\mathcal D_X)
 \ \simeq\ D^b_c(X^{\mathrm{an}},\mathbb C).
\]
左边的对象是左 \(\mathcal D_X\) 模的有界复形，其上同调模正则且 holonomic；右边的对象是复向量空间层的有界复形，其上同调层沿有限代数分层局部常值且纤维有限维。正则性包括无穷远处的正则条件。该函子把单个正则 holonomic 左 \(\mathcal D_X\) 模送到反常层。
\end{theorem}
\symidx{\(\operatorname{DR}_X(M)\)：左 D-模的移位 de Rham 复形}
\symidx{\(D^b_{\mathrm{rh}}(\mathcal D_X),D^b_c(X^{\mathrm{an}},\mathbb C)\)：正则 holonomic D-模与构造性层的有界导出范畴}
\begin{proofsketch}
先看正规交叉边界上的正则亚纯联络。它的 de Rham 复形在局部坐标中是由互相交换的
\(\partial_i-A_i/z_i\) 组成的 Koszul 复形。逐坐标 Laurent 展开将其上同调与穿孔图上局部系统的上同调比较；正则性保证以亚纯函数或以所有解析函数计算时，包含映射给出同一个上同调。这是证明与开嵌入的延拓相容所需的局部比较。对一般边界，先消解成正规交叉边界，再用适当直接像下降；该比较见\cite[Theorem 5.2.24 and Remark 5.2.25, pp. 151--153]{HottaTakeuchiTanisaki2008DModules}。

全忠实的验证使用对偶、外张量积、对角线逆像和到点的直接像：两个对象的导出 Hom 可由这些操作表达。正则增长的局部比较使 de Rham 函子与上述操作相容，因而把该表达转成两个可构造复形的导出 Hom。取零次上同调便得到 Hom 集合的双射，各次上同调则同时控制扩张。这里还须核对所得同构与 de Rham 实际诱导的同态一致，而不能仅比较两个复形的抽象同构类。

再证本质满射。可构造导出范畴由
\(Ri_*\mathcal L\) 生成，其中 \(i:S\hookrightarrow X\) 是光滑局部闭代数子簇的仿射嵌入，\(\mathcal L\) 为其有限秩局部系统。对有限分层使用开闭分解给出的好三角并逐层归纳，就得到这一生成陈述。由\cref{b5:thm:D-rh-connections}，在 \(S\) 上取对应的正则联络 \(N\)；微分算子直接像（记为 \(i_+\)，以转移微分算子双模及导出张量积定义）
\(i_+N[-\dim S]\) 是正则 holonomic 复形，其 de Rham 复形为 \(Ri_*\mathcal L\)。闭嵌入的法向部分可由上一节的
\(\mathbb C[\partial_1,\ldots,\partial_c]\otimes N\)
模型理解，开嵌入则使用最初的正则亚纯比较。函子的像已包含这些生成元，全忠实又允许将所有连接映射与映射锥提升，故它覆盖整个可构造导出范畴。

这里省略正则 holonomic 模在这些操作下的保持性、de Rham 与各操作相容的比较，以及导出 Hom 同构的自然性检查；这些是一般对应的大部分技术内容。相关证明及自然性问题见\cite[Sections 7.1--7.2, Theorems 7.2.2 and 7.2.5, Remark 7.2.3, pp. 174--177]{HottaTakeuchiTanisaki2008DModules}。在纯 \(n\) 维光滑簇上，一个联络的 de Rham 复形是相应局部系统移位 \([n]\)。若在一个光滑 \(d\) 维闭子簇上取联络，再作闭嵌入直接像，则得到该子簇上局部系统移位 \([d]\) 的直接像。一般正则 holonomic 模在各分层片上的限制未必只集中在一个次数；反常性须由支撑与余支撑条件判断，其中使用分层片自身的维数。一般非正则模不满足前述增长比较，因而不属于这个对应。
\end{proofsketch}

在一维情形，可以直接看见为何单个 D-模的像也需要是复形。对 \(X=\mathbb A^1_{\mathbb C}\)，解析 Poincaré 引理说明
\[
\operatorname{DR}(\mathcal O_X)
=[\mathcal O_{X^{\mathrm{an}}}\xrightarrow{\,\mathrm{d}\,}
\Omega^1_{X^{\mathrm{an}}}][1]
\simeq\mathbb C_{X^{\mathrm{an}}}[1].
\]
未移位的复形只有零次上同调常数层，移位后它位于次数 \(-1\)。另一方面，对 \(M=\mathcal D_X/\mathcal D_Xx\)，上一节的法向模型为 \(\mathbb C[\partial]\)，其 de Rham 微分是乘以 \(\partial\)。这个算子单射，余核由常数项生成，且在 \(x\ne0\) 处整个模为零，故
\[
\operatorname{DR}(\mathcal D_X/\mathcal D_Xx)
\simeq\mathbb C_{\{0\}}[0].
\]
两个普通 D-模的支撑维数分别为一与零，而像的非零上同调分别位于 \(-1\) 与 \(0\)。这里的 Verdier 对偶是相对于对偶化复形取导出层 Hom；在光滑复 \(n\) 维空间上，对偶化复形为 \(\mathbb C_X[2n]\)。反常层的支撑条件要求 \(\dim\operatorname{Supp}\mathcal H^j\le-j\)，对其 Verdier 对偶还须满足同样条件；这两个模型恰展示了维数与次数的配合。若只用集中在同一次数的普通层，就不能同时包含这两种对象。

\ProseParagraphBreak
正维真支撑的例子也可以直接算出。取 \(X=\mathbb A^2_{\mathbb C}\)、\(S=\{x=0\}\)，以及
\[
 M=D_2/(D_2x+D_2\partial_y),\qquad
 i:S\hookrightarrow X.
\]
把 \(x\) 与 \(\partial_y\) 移到右端，正规表达给出 \(M\simeq\mathbb C[\partial_x]\otimes\mathbb C[y]\)：\(x\) 按 \(-\mathrm{d}/\mathrm{d}\partial_x\) 作用，\(\partial_y\) 对第二因子求导。因此商滤过的特征簇为 \(V(x,\xi_y)\)，维数为 \(2\)，所以它是 holonomic 模。解析化后，未移位 de Rham 复形是法向两项复形与切向复形的总化。前者的微分为乘以 \(\partial_x\)，核为零、余核在次数 \(1\) 留下常数；后者在 \(S\) 上的微分为通常求导，由解析 Poincaré 引理只在次数 \(0\) 留下常数层。总复形于是只在次数 \(1\) 有上同调 \(i_*\mathbb C_S\)，再作环境维数移位 \([2]\)，得到
\[
 \operatorname{DR}_X(M)\simeq i_*\mathbb C_S[1].
\]
在二维环境中，平凡联络、这条曲线上的平凡联络直接像、点支撑模的像分别带有移位 \([2]\)、\([1]\)、\([0]\)。这些是光滑支撑上的具体模型，并不意味着任意同维支撑的 holonomic 模都有同样简单的像。

\subsection{与分析学、几何表示论的联系}
\label{b5:subsec:D02:04}

下面计算导出 Hom，比较两种系数环所产生的上同调。令 \(D=A_1(\mathbb C)\)，\(M=D/D\partial\)。右乘 \(\partial\) 是左 \(D\) 线性单射，因为正规表达保证 \(D\) 无零因子，故有自由消解
\[
 0\longrightarrow D\xrightarrow{\,R_\partial\,}D\longrightarrow M\longrightarrow0.
\]
对左模 \(\mathbb C[x]\) 取 \(\operatorname{Hom}_D\bigl(-,\mathbb C[x]\bigr)\)，所得两项复形为
\[
 \mathbb C[x]\xrightarrow{\mathrm{d}/\mathrm{d}x}\mathbb C[x].
\]
其零次上同调为常数，一次上同调为零，因为多项式在特征零下总可逐项积分。

\ProseParagraphBreak
若改在穿孔直线上，用
\(D^\times=\mathbb C[x,x^{-1}][\partial;\mathrm{d}/\mathrm{d}x]\) 及同样的常数方程模，计算变为
\[
 \mathbb C[x,x^{-1}]\xrightarrow{\mathrm{d}/\mathrm{d}x}\mathbb C[x,x^{-1}].
\]
核仍为 \(\mathbb C\)，余核却是一维，由 \(x^{-1}\) 的类生成：除 \(x^{-1}\) 外，每个 Laurent 单项式都有 Laurent 原函数。这个余核记录原函数无法在穿孔直线全局选取的障碍，与沿圈积分 \(\int \mathrm{d}x/x=2\pi i\) 相呼应。此处算的是全局代数模的导出 Hom；局部解析求导仍满射，所以不能把这项全局上同调误称为每个点的局部解障碍。

\begin{proposition}\label{b5:prop:D-laurent-monodromy-cohomology}
设 \(\lambda\in\mathbb C\)，\(R=\mathbb C[x,x^{-1}]\)，并设
\(T_\lambda:R\to R\)、\(T_\lambda(f)=xf'-\lambda f\)。在 \(X=\mathbb C^\times\) 上，记 \(\mathcal L_\lambda\) 为解联络 \(\nabla=\mathrm{d}-\lambda\cdot\mathrm{d}x/x\) 的水平截面局部系统，其正向单值为 \(T=e^{2\pi i\lambda}\)。则
\[
 \begin{aligned}
 \dim_{\mathbb C}\ker T_\lambda
 &=\dim_{\mathbb C}\operatorname{coker}T_\lambda\\
 &=\dim_{\mathbb C}H^0(X,\mathcal L_\lambda)
  =\dim_{\mathbb C}H^1(X,\mathcal L_\lambda)\\
 &=\begin{cases}1,&\lambda\in\mathbb Z,\\0,&\lambda\notin\mathbb Z.\end{cases}
 \end{aligned}
\]
整数情形下，核由 \(x^\lambda\) 生成，余核由 \(x^\lambda\) 的类生成。
\end{proposition}
\begin{proof}
Laurent 单项式是 \(R\) 的基，且
\(T_\lambda(x^m)=(m-\lambda)x^m\)。若 \(\lambda\notin\mathbb Z\)，对有限 Laurent 和逐项除以非零系数 \(m-\lambda\) 即得逆映射。若 \(\lambda\in\mathbb Z\)，只有 \(x^\lambda\) 这一基向量被零化，而像恰为缺少这一项的 Laurent 多项式，得到所述核与余核。

拓扑方面，\(\mathbb C^\times\) 形变收缩到圆周。用两个可缩开弧覆盖圆周，其交集有两个可缩分量；局部系统在其中一个分量上的过渡为恒等，在另一个为 \(T\)。这份零调覆盖给出的 Čech 微分可以写为
\[
 \mathbb C^2\longrightarrow\mathbb C^2,\qquad
 (a,b)\longmapsto(b-a,b-Ta).
\]
先令 \(c=b-a\)，再将第二输出减去第一输出，便拆成一个恒等映射和 \(a\mapsto(1-T)a\)。因此上同调由两项复形 \(\mathbb C\xrightarrow{T-1}\mathbb C\) 计算，分别为其核与余核。\(T=1\) 当且仅当 \(\lambda\in\mathbb Z\)，得出相同的维数。

在代数形式中，\(\nabla f=(f'-\lambda f/x)\,\mathrm{d}x\)；用基 \(\mathrm{d}x/x\) 表示一形式，微分正是 \(T_\lambda\)。整数情形下换水平基 \(x^\lambda\)，复形变为通常求导的复形，而余核代表 \(x^\lambda\,\mathrm{d}x/x\) 化为 \(\mathrm{d}x/x\)，沿正向圆周积分为 \(2\pi i\)。它确实给出非零的拓扑上同调类；非整数情形两边都为零。这个算例直接核对了代数微分复形与局部系统上同调的比较。
\end{proof}

同一个参数也能检验解函子与 de Rham 函子的对偶方向。令
\(Q_\lambda=D^\times/D^\times(x\partial-\lambda)\)，以 \(u\) 为其循环生成元，则 \(\partial u=\lambda u/x\)。到普通函数模的同态由 \(u\mapsto y\) 决定，关系要求 \(y'=\lambda y/x\)，所以在单连通穿孔区域内
\[
 \operatorname{Hom}_{D^\times}(Q_\lambda,\mathcal O)
 =\mathbb Cx^\lambda.
\]
这正是上一命题的解局部系统；自由消解所得导出 Hom 的微分为 \(T_\lambda\)。但模自身的截面 \(fu\) 水平，要求的是
\[
 \partial(fu)=\left(f'+\frac\lambda x f\right)u=0,
 \qquad f=cx^{-\lambda}.
\]
因此 \(Q_\lambda\) 自身的水平局部系统具有单值 \(e^{-2\pi i\lambda}\)，解局部系统则具有 \(e^{2\pi i\lambda}\)。例如 \(\lambda=1/3\) 时两者互为逆且不同。\(D\) 线性同态 \(Q_\lambda\to\mathcal O\) 就是 \(Q_\lambda^*\) 的水平截面；相反，\(\operatorname{DR}(Q_\lambda)\) 使用 \(Q_\lambda\) 自身的联络。其移位 \([1]\) 只改变复形次数，不会把这两种单值互换。附录 E 采用 \(\nabla e=-\lambda e/x\) 来写系统 \(y'=\lambda y/x\)，因此对应的是这里的对偶微分模 \(Q_\lambda^*\)：其对偶基元满足 \(\partial u^*=-\lambda u^*/x\)，水平截面 \(fu^*\) 满足 \(f'=\lambda f/x\)，单值正是 \(e^{2\pi i\lambda}\)。

\ProseParagraphBreak
几何表示论还把 Lie 代数的包络代数模与旗簇上的适当扭曲 D-模联系起来。局部化定理需要指定中心特征、参数与正则性条件。本附录介绍了微分算子、特征簇和层复形等所需概念；局部化与表示分类的系统理论可参阅
\cite[Chapter 11]{HottaTakeuchiTanisaki2008DModules}。


\begin{chaptersummary}
Leibniz 法则成为 Weyl 代数的交换关系，线性方程则成为自由左模中的关系。阶滤过把这种非交换结构转为余切空间上的交换代数；特征簇据此记录约束的几何形状。它不能用任意滤过的符号零点集代替，计算一般左理想时也必须掌握整个分次理想。

\ProseParagraphBreak
形式级数递推、解析延拓和扇形渐近各自保存不同的信息。正则联络可以由局部系统描述，但非正则方程还可能带有 Stokes 数据；全局原函数的障碍又需要上同调记录。因此 Riemann--Hilbert 对应中的正则性、构造性和导出范畴都有具体作用。
\end{chaptersummary}

\BookExercises

\begin{exercise}
设 \(k\) 为特征零域，\(m\ge1\)。在 \(A_1(k)\) 中求 \(\partial^m x\) 与 \(\partial^2x^3\) 的正规表达。
\end{exercise}
\begin{solution}
由 \(\partial x=x\partial+1\) 归纳得
\(\partial^m x=x\partial^m+m\partial^{m-1}\)。先算
\(\partial x^3=x^3\partial+3x^2\)，再求导一次，得
\[
 \partial^2x^3=x^3\partial^2+6x^2\partial+6x.
\]
额外项来自求导对系数的作用，不能按交换多项式相乘。
\end{solution}

\begin{exercise}
在特征零域 \(k\) 上证明 Weyl 代数没有非零有限维左模。说明这不妨碍存在在 \(k(x)\) 上有限维的微分模。
\end{exercise}
\begin{solution}
若模在 \(k\) 上维数为正整数 \(d\)，对作用关系
\([\partial,x]=I\) 取矩阵迹，左边为零，右边为 \(d\ne0\)，矛盾。附录 E 的\cref{b5:def:E-differential-module}所定义的微分模，是有限维 \(k(x)\) 向量空间连同满足 Leibniz 法则的联络算子；它不同于记录导子的 Kähler 微分模。若这样的微分模非零，任取非零向量 \(v\)，若各个 \(x^jv\)、\(j\ge0\) 在 \(k\) 上线性相关，就有非零 \(P(x)\in k[x]\) 满足 \(P(x)v=0\)。但 \(P(x)\) 在 \(k(x)\) 中可逆，这会推出 \(v=0\)，矛盾。因此这些向量线性无关，底层 \(k\) 向量空间必无限维，上述有限矩阵的迹论证并不适用。
\end{solution}

\begin{exercise}
设 \(D=A_1(\mathbb C)\)。构造一个左 \(D\) 模，使其到任意函数模 \(V\) 的同态对应方程组 \(y_1'=y_2\)、\(y_2'=-y_1\) 的解。把它化为单个循环模。
\end{exercise}
\begin{solution}
取 \(D^2\) 的基 \(e_1,e_2\)，模去关系
\(\partial e_1-e_2\)、\(\partial e_2+e_1\) 生成的左子模。同态把基送到 \(y_1,y_2\)，故关系恰好成为所给方程。由第一关系 \(e_2=\partial e_1\)，第二关系化为
\((\partial^2+1)e_1=0\)。映射 \(D/D(\partial^2+1)\to M\)、\(1\mapsto e_1\) 与
\(e_1\mapsto\overline1,e_2\mapsto\partial\overline1\) 互逆。
\end{solution}

\begin{exercise}
在 \(D=A_1(\mathbb C)\) 中，求 \(P=x^2\partial^3+\partial-x^5\) 的阶主符号与 Bernstein 主符号，说明相应零点集为何不同。
\end{exercise}
\begin{solution}
最大微分阶为 \(3\)，故阶主符号是 \(x^2\xi^3\)。Bernstein 次数把 \(x,\partial\) 都计为一次；首项和末项的次数均为 \(5\)，故符号是
\(x^2\xi^3-x^5=x^2(\xi^3-x^3)\)。两个零点集分别为
\[
V(x)\cup V(\xi),\qquad
V(x)\cup\bigcup_{\zeta^3=1}V(\xi-\zeta x).
\]
前者是两条坐标轴之并，后者是 \(x=0\) 轴与三条 \(\xi=\zeta x\) 直线之并；按本附录的定义，前者是阶滤过给出的特征簇；后者是 Bernstein 滤过下的符号零点集，不能用它替代前者。
\end{solution}

\begin{exercise}
在 \(D=A_1(\mathbb C)\) 中，令 \(I=D\partial^2+Dx^2\)。证明 \(I=D\)，并说明为何两条生成关系的主符号仍不足以给出特征簇。这里不能直接把 \(x^2v=\partial^2v=0\) 换成 \(xv=\partial v=0\)。
\end{exercise}
\begin{solution}
令 \(E=x\partial\)。由 Weyl 关系可算出
\[
 A=\partial^2x^2=E^2+3E+2\in I,\qquad
 B=x^2\partial^2=E^2-E\in I.
\]
所以 \(A-B=4E+2\in I\)，再利用左理想对左乘封闭，得到
\(4B-E(A-B)=-6E\in I\)。特征零下可除以 \(-6\)，故 \(E\in I\)，随后 \(2\in I\)、\(1\in I\)，于是 \(I=D\)。商模为零，特征簇为空；两个生成元的阶主符号却为 \(\xi^2,x^2\)，共同零点仍是 \((0,0)\)。这些符号没有记录从算子组合中得到的单位关系。
\end{solution}

\begin{exercise}
对 \(D=A_1(\mathbb C)\) 及 \(a\in\mathbb C\)，计算
\(D/D\bigl((x-a)\partial\bigr)\) 的特征簇，并求该方程在 \(\mathbb C[x]\) 中的解。
\end{exercise}
\begin{solution}
\cref{b5:prop:D-principal-characteristic}给出
\(\operatorname{Ch}=V\bigl((x-a)\xi\bigr)\)，即零截面与 \(a\) 上余切纤维之并。多项式解满足 \((x-a)y'=0\)，多项式环无零因子，故 \(y'=0\)，解仅为常数。特征簇的额外分支不等于新增的多项式解，它记录方程在特殊点的退化。
\end{solution}

\begin{exercise}
在 \(\mathbb A^2_\mathbb C\) 的平凡线丛上设
\(\nabla_x=\partial_x+a(x,y)\)、\(\nabla_y=\partial_y+b(x,y)\)，其中 \(a,b\in\mathbb C[x,y]\)。
求成为 D-模的条件，并判定 \(a=y,b=x\) 与 \(a=y,b=0\) 两例。
\end{exercise}
\begin{solution}
对任意函数作用后相减，所有一阶求导项抵消，留下乘法算子
\(\partial_xb-\partial_ya\)。所以条件为 \(\partial_xb=\partial_ya\)。
第一例两边都是 \(1\)，成立；第二例两边为 \(0,1\)，不成立。Leibniz 法则自身不足以保证多个微分方向相容。
\end{solution}

\begin{exercise}
在 \(\mathbb C[\![x]\!]\) 中，对形式方程 \(y'=xy\)、\(y(0)=1\)，求全部形式系数，并决定是否收敛。
\end{exercise}
\begin{solution}
设 \(y=\sum a_mx^m\)，有 \(a_1=0\) 及
\((m+1)a_{m+1}=a_{m-1}\) 对 \(m\ge1\) 成立。于是奇数系数为零，\(a_{2j}=1/(2^jj!)\)，即
\(y=\exp(x^2/2)\)。比值判别表明对每个 \(x\) 都收敛；这里收敛性是另外核验的，不能只由形式递推推出。
\end{solution}

\begin{exercise}
设 \(\lambda\in\mathbb C\)、\(m\in\mathbb Z\)，令
\(N=\begin{pmatrix}0&1\\0&0\end{pmatrix}\)、\(A=\lambda I+N\)。在穿孔复圆盘上求系统 \(Y'=AY/z\) 的基本解矩阵与单值。求标量换基 \(Y=z^mV\) 后的系统，并说明整数指数平移是否能够消去对数项。
\end{exercise}
\begin{solution}
在通用覆盖上取 \(Z=z^\lambda(I+N\log z)\)。由 \(N^2=0\) 直接求导得 \(Z'=AZ/z\)，且 \(\det Z=z^{2\lambda}\ne0\)。正向延拓后，
\[
 Z\longmapsto ZT,\qquad
 T=e^{2\pi i\lambda}(I+2\pi iN).
\]
换基给出 \(V'=(A-mI)V/z\)，指数变为 \(\lambda-m\)，单值仍为 \(T\)。它只改变标量指数，没有消去 \(N\log z\)；单值亚纯换基沿闭路返回同一矩阵，至多使基点表示中的单值矩阵共轭；非恒等的幺幂矩阵不可能经共轭变为恒等，所以这种换基也不能消去对数单值。对数项不能只靠整数指数平移消失。
\end{solution}

\begin{exercise}
判断秩一方程 \(y'=(z^{-3}+2z^{-1})y\) 在 \(0\) 的正则性，求一个解析基本解和单值。解释单值能否判定正则性。
\end{exercise}
\begin{solution}
积分得 \(y=z^2\exp\bigl(-1/(2z^2)\bigr)\)，它在穿孔圆盘单值，故单值为 \(1\)。但是系数的三阶极点不能由亚纯 \(g'/g\) 消去，由定义~\ref{b5:def:D-regular-singular}前后的秩一判别，它非正则。单值为 \(1\) 并不意味着正则或解析可去。
\end{solution}

\begin{exercise}
在任意闭扇形 \(|\arg x|\le\theta<\pi/2\) 中证明 \(e^{-1/x}=O(|x|^N)\) 对每个整数 \(N\ge0\) 成立。由此说明发散形式解的渐近展开不能唯一确定例~\ref{b5:exm:D-euler-stokes}的解析解。
\end{exercise}
\begin{solution}
有 \(\operatorname{Re}(1/x)\ge\cos\theta/|x|\)，故
\(|e^{-1/x}|\le e^{-\cos\theta/|x|}\)。令 \(s=1/|x|\)，函数
\(s^Ne^{-\cos\theta s}\) 有界，给出所需估计。因此给某个扇形解加上 \(ce^{-1/x}\)，所有有限阶渐近系数都不改变，却一般得到不同解析函数。
\end{solution}

\begin{exercise}
设 \(R=\mathbb C[x,x^{-1}]\)、\(A=\lambda I+N\)，其中 \(\lambda\in\mathbb C\)、\(N=\begin{pmatrix}0&1\\0&0\end{pmatrix}\)。求 \(T_A:R^2\to R^2\)、\(T_A(f)=xf'-Af\) 的核与余核，并与单值 \(T=e^{2\pi i\lambda}(I+2\pi iN)\) 的 \(T-I\) 相比较。
\end{exercise}
\begin{solution}
对 \(c\in\mathbb C^2\)，有
\(T_A(x^mc)=x^m\bigl((m-\lambda)I-N\bigr)c\)。当 \(m\ne\lambda\) 时，该矩阵的逆为
\[
 \frac1{m-\lambda}I+\frac1{(m-\lambda)^2}N,
\]
故这些 Laurent 次数均不贡献核或余核。若 \(\lambda\notin\mathbb Z\)，所有次数都可逆，两者为零。若 \(\lambda\in\mathbb Z\)，唯一剩下的次数为 \(m=\lambda\)，其矩阵为 \(-N\)；核由 \(x^\lambda e_1\) 生成，余核由 \(x^\lambda e_2\) 的类生成，两者均为一维。拓扑单值在非整数情形下使 \(T-I\) 可逆；整数情形则有 \(T-I=2\pi iN\)，其核为 \(\mathbb Ce_1\)，余核由 \(e_2\) 的类生成。若 \(N=0\)，整数情形下的核与余核各有维数二；这里非零的 Jordan 块使两者各降为一维，对应非平凡对数单值。
\end{solution}
